\section{Conclusion and Visionary Future Directions}
\label{sec:conclusion}

In this work, we challenged the foundational assumption of deep learning that neural network parameter updates must be merged within static, real-valued Euclidean coordinates. Instead, we introduced \textbf{PhaseMerge}, a continuous wave superposition framework that projects task-specific parameter updates into complex-valued Fourier frequency-space. By modeling fine-tuning updates as continuous wavefunctions, we enabled a differentiable phase-rotation mechanism.

Our empirical results on Vision Transformers across highly conflicting vision tasks confirm that PhaseMerge is highly competitive and robust, outperforming the traditional static baseline (Uniform TA) and static frequency-filtering baselines (FREE-Merging) across FP32 and PTQ regimes, and matching the performance of strong, constrained real-space baselines under extreme quantization. Crucially, by parameterizing phase-shifts as a uniform rotation per layer ($r=1$), Uniform PhaseMerge (U-PhaseMerge) introduces an implicit matrix-basis regularizer on the parameter updates. This successfully avoids catastrophic overfitting to validation noise under moderate calibration data scarcity ($M=16$). Furthermore, because our phase-shift profile is smooth and operates in Fourier space, it rejects localized rounding discretization thresholds, ensuring excellent robustness across post-training quantization schema shifts.

\textbf{Visionary Future Frontiers.} We believe that PhaseMerge represents only the first step in a larger, wave-theoretic paradigm shift for neural representation and model merging. We propose several concrete and exciting directions for future research:
\begin{itemize}
    \item \textbf{Input-Dependent Phase Modulation:} Instead of optimizing a static phase grid offline, future work can explore real-time, input-dependent wave-front modulating networks. This would allow weight matrices to dynamically adjust their phase profiles on the fly to constructively align with the specific frequency requirements of each incoming sample.
    \item \textbf{Phase-Decay and Coordinate Regularization:} Investigating soft $L_2$ phase regularization penalties (such as phase decay) during optimization to prevent phase angles from drifting excessively from their initial coordinates under larger calibration sizes ($M \ge 32$).
    \item \textbf{Extension to Spatial Convolutional Kernels:} While dense weights lack an innate physical grid correlation, convolutional kernels (such as those in CNNs, ResNets, or ConvNeXTs) possess actual, physical height and width dimensions with strong spatial correlation. Applying PhaseMerge to convolutional filters would resolve any abstract neuron-order concerns, aligning the 2D/3D Fast Fourier Transform perfectly with the physical topology of the network.
    \item \textbf{Multi-Dimensional Spectral Fusion:} Extending Fourier analysis and continuous phase synchronization to multi-dimensional spectral representations (e.g., 3D or tensor-based FFTs) for higher-order convolutional and attention layers.
    \item \textbf{Scaling to Generative Foundations:} Applying PhaseMerge to multi-billion parameter Large Language Models (LLMs) and Diffusion experts to achieve zero-overhead, multi-modal parameter fusion directly on resource-constrained edge hardware.
\end{itemize}
By moving beyond static coordinate compromises, we open up a vast, mathematically rich field where parameter optimization and wave mechanics converge.
